\begin{longtable}{@{}>{\raggedright\arraybackslash}p{2.4cm}rr>{\raggedright\arraybackslash}p{1.3cm}rr>{\raggedright\arraybackslash}p{4.8cm}@{}}
\caption{Inputs that matter for the results, ranked by their effect on any of the three screened ships: the largest change in the cost saving of a rig between the low and high bound, and the largest change in the cost of a ship without a rig at either bound.}
\label{tab:inputs}\\
\toprule
Input & Value & Range & Unit & Saving & Cost & Source \\
 & & & & [pt] & [\%] & \\
\midrule
\endfirsthead
\multicolumn{7}{@{}l}{\cref{tab:inputs} (continued)}\\
\toprule
Input & Value & Range & Unit & Saving & Cost & Source \\
 & & & & [pt] & [\%] & \\
\midrule
\endhead
\bottomrule
\endfoot
rotor cost per area & \num{2040} & \num{1260}--\num{4270} & USD/m\textsuperscript{2} & 14.0 & 0.0 & Manufacturing cost per m\textsuperscript{2} of height times diameter, the slope of the same line. \\
wing cost per area & \num{665} & \num{310}--\num{1760} & USD/m\textsuperscript{2} & 11.9 & 0.0 & Manufacturing cost per m\textsuperscript{2} of sail from 0.13 t/m\textsuperscript{2} (EMSA 2023, Table 4; OceanWings; Oceanbird) of tubular and stiffened steel, composite and slewing parts at NREL and WISDEM costs. \\
canopy making cost per area & \num{42} & \num{16}--\num{179} & USD/m\textsuperscript{2} & 10.5 & 0.0 & Goal for the manufacturing cost per m\textsuperscript{2} of projected area of a ram-air canopy besides its cloth: per flat m\textsuperscript{2} 1.0 h of sewing at 26 USD/h, 7 USD of bridle lines and 3 USD of reinforcements, tapes and hardware, 1.17 flat m\textsuperscript{2} per projected m\textsuperscript{2}. Today 1.5-2.25 h per flat m\textsuperscript{2}, 50 h of manufacturing per paraglider (SKY Paragliders 2017) over 22-33 m\textsuperscript{2}; the US G-12 cargo parachute implies 0.6-1.0 h at series volume if materials take 57 \% of sales as in 2000, up to 1.75 h at the 25-50 \% of other years (BIS 2004). \\
e ammonia price & \num{47} & \num{30}--\num{82} & USD/GJ & 10.5 & 23.5 & E-ammonia bunkered in 2030: 47-82 USD/GJ (Lloyd's Register \& UMAS 2020, Table B3, p. 43); 30-43 USD/GJ (MMMCZCS 2021, Fuel Option Scenarios, p. 12), the low bound. The base is the lower of the first, 0.64 of e-methanol's 73 there. \\
kite max force per area & \num{0.5} & \num{0.4}--\num{1} & kN/m\textsuperscript{2} & 10.3 & 0.0 & Maximum wing loading on the projected area. 450 N/m\textsuperscript{2} (Faggiani \& Schmehl 2018); Kitepower Falcon 3.5 t on 60-80 m\textsuperscript{2} (0.43-0.57 kN/m\textsuperscript{2}). SkySails design load on flat area: 160 kN on 225 m\textsuperscript{2}, about 0.7 kN/m\textsuperscript{2} (Paulig et al. 2013), 0.9 on the projected area at 1.28 flat over projected (Folkersma et al. 2020). The high bound is the wing load of up to 1 kN/m\textsuperscript{2} in SkySails' operating experience (SkyPower100 final report, Junge et al. 2023, sec. I.4). The canopy cloth, seams and reinforcements follow the limit (canopy\_cloth\_force\_exponent), so the bounds carry them too. \\
wind shear exponent & \num{0.05} & \num{0.04}--\num{0.08} & -- & 10.1 & 0.0 & Exponent of the mean power-law wind profile over the sea from 10 m to kite height: matches the mean ERA5 model-level profiles over the study region in 2018 (1.12, 1.18 and 1.22 times the 10 m wind at 100, 270 and 540 m; own analysis of Hersbach et al. 2020). The low bound is the trade-wind belt (1.17 at 540 m), the high bound the North Atlantic north of 40 deg N and the IJmuiden lidar, which follow a logarithmic profile with 0.2 mm roughness up to 315 m (Kalverla 2019, Ch. 2; Schelbergen et al. 2020). \\
canopy life & \num{3000} & \num{1200}--\num{6900} & flight hours & 10.0 & 0.0 & Development goal for the life of a canopy of 200 g/m\textsuperscript{2} cloth: what development can reach, not where things stand today. No life is measured on a kite. Reachable and on the careful side of what makers state: below the 3,500-4,000 h of their longest statements and the 3,800 h that strength decay under UV and a wear rate from makers' statements gives. Each decay input changed alone gives the bounds. \\
kite launch max wave height & \num{3} & \num{2}--\num{4} & m & 9.5 & 0.0 & Estimate, no source gives a number: highest significant wave height at which a kite is launched, one metre below the stay-up limit because a launch is the step most exposed to ship motion. The SkySails supplier reported limited launch opportunities in heavy sea in the first test phase (CE Delft 2016). \\
kite max wave height & \num{5} & \num{4}--\num{7} & m & 8.5 & 0.0 & Estimate, no source gives a number: highest significant wave height at which a kite aloft stays up. A kite far above the sea feels the ship's motion only through its tether, so we set this one metre above the landing limit. \\
kite min turn radius & \num{2.3} & \num{1.5}--\num{4.3} & -- & 8.4 & 0.0 & Smallest radius of the kite's path through the air over its projected span: the SkySails turn rate law gives a yaw rate g v\_a delta, so a turn radius 1/(g delta), with g = 0.03-0.05 rad/m for the kite of 160 m\textsuperscript{2} flat, 125 m\textsuperscript{2} projected, and curve deflections delta of 40-70 \% of the steering range (Erhard \& Strauch 2013, Eq. 7, Section VIII); 1/(0.04 x 0.55) = 45 m, with bounds 29 m and 83 m, over the 19.4 m span of that area at aspect ratio 3. \\
vlsfo price & \num{16} & \num{12}--\num{19} & USD/GJ & 7.4 & 9.0 & Rotterdam VLSFO averaged \textasciitilde{}663 USD/t over 2026 and was 778 USD/t on 29 Oct 2026 (Ship \& Bunker), at 41 GJ/t. \\
discount rate & \num{0.08} & \num{0.06}--\num{0.12} & -- & 6.5 & 23.6 & Damodaran (Jan 2026) puts the WACC of Shipbuilding \& Marine at 6.7 \%; a small private owner pays more, a startup up to 12 \%. \\
carbon coverage & \num{0} & \num{0}--\num{1} & -- & 6.3 & 14.2 & A 3,500 dwt coaster measures \textasciitilde{}2,500 GT (Damen CF3850: 2,570 GT), below the 5,000 GT threshold of EU ETS and FuelEU, so it pays nothing today. A July 2026 Commission proposal would bring smaller ships into the ETS from 2031; Rotterdam-Guadeloupe is then an intra-EU voyage counted in full. \\
market price factor & \num{2} & \num{1.6}--\num{2.5} & -- & 4.9 & 0.0 & Market price over manufacturing cost for wings, rotors and kites alike: production overhead of 1.37 times direct cost (NREL blade cost model) over one minus a gross margin of 30 \%, between Vestas (13 \%) and European machinery (45 \%, Damodaran 2026). \\
tether stress & \num{4.0e8} & \num{3.0e8}--\num{4.8e8} & Pa & 4.9 & 0.0 & Working stress on the nominal diameter of the rope at design factor 3: 12-strand HMPE, 880 kN mean breaking load at 30.6 mm nominal, 399 MPa (Bosman et al. 2013, Table 33.4), the starting factor Bosman et al. and Fechner \& Schmehl (2013) advise. Also \textasciitilde{}0.40 GPa in Faggiani \& Schmehl (2018). The low bound is design factor 4, the high bound 40 \% of break load, above which the rope deteriorates rapidly in bending (Bosman et al. 2013). \\
rotor max wind & \num{20} & \num{15}--\num{25} & m/s & 4.9 & 0.0 & Apparent wind above which the rotors stop spinning and only their drag remains (Tillig \& Ringsberg 2020, Table 3). \\
newbuild price & \num{1.19e7} & \num{1.0e7}--\num{2.0e7} & USD & 2.9 & 19.3 & Price of a conventional cargo ship of the reference deadweight from an Asian yard: least-squares power law through Udupi Cochin's 3,800 dwt general cargo ships (\textasciitilde{}USD 11.6 M each, Onmanorama 15 Jun 2023) and the KOBC 33k/64k/82k dwt bulker assessments of USD 31.25/35.78/38.22 M (ShipUniverse, 8 Oct 2026). A Dutch build costs more: Royal Bodewes 6,750 dwt EcoTrader \textasciitilde{}EUR 17 M (Meriaura, Apr 2024). \\
tether rope density & \num{830} & \num{590}--\num{1130} & kg/m\textsuperscript{3} & 4.6 & 0.0 & Rope mass over its nominal cross-section: 60.8 kg per 100 m at 30.6 mm, 827 kg/m\textsuperscript{3} (Bosman et al. 2013, Table 33.4, coated 12-strand). Their worked-in diameter, 26.0 mm, would give 1,145 kg/m\textsuperscript{3}, above solid HMPE at 970 kg/m\textsuperscript{3}, so it cannot be the section that meets the wind. Low: Marlow D12 78 uncoated, 4.65 kg per 100 m at 10 mm. High: the SkySails tether with its power and data cable, 0.9 kg/m for 320 kN (Paulig et al. 2013). \\
cargo per teu & \num{10.2} & \num{8}--\num{14} & t & 3.2 & 17.9 & Mean cargo mass per loaded TEU handled in EU ports (Eurostat). \\
thrust deduction ratio & \num{0.6} & \num{0.5}--\num{0.75} & -- & 4.4 & 4.0 & t = 0.6 w (MAN ES 2018, Taylor), which equals MAN's alternative t = 0.27 C\_B at C\_B 0.83. \\
canopy cloth base price & \num{6} & \num{2}--\num{17} & USD/m\textsuperscript{2} & 4.3 & 0.0 & Goal for the price of a sewn m\textsuperscript{2} of cloth, with 15 \% offcuts, at no cloth weight: the retail line through polyester and nylon kite cloth and Dacron sailcloth of 29-385 g/m\textsuperscript{2} ex VAT, 11 USD/m\textsuperscript{2} plus 0.031 USD/m\textsuperscript{2} per g/m\textsuperscript{2} (extremtextil.de, Sailrite, Jones Airfoils, 2026-11-01), scaled to USD 9.4 at 200 g/m\textsuperscript{2}, what a maker buying by the roll and from the mill can reach. No mill quote for this cloth. \\
rig maintenance fraction & \num{0.02} & \num{0.01}--\num{0.05} & 1/year & 3.9 & 0.0 & Maintenance and repair of wings, rotors and kite hardware over their price: 2 \% a year for every rig but kites, and 5 \% for kites with their wear (EMSA 2023, p.45 and Table 7), the high bound, as the canopy and tether are charged by wear here; Norsepower's figure for rotors is also 2 \% (Tillig et al. 2020, p.15). The low bound is an estimate. \\
kite flight drag & \num{0.08} & \num{0}--\num{0.12} & -- & 3.9 & 0.0 & Drag on the projected area that the design-shape model lacks in flight: deformation, control unit and steering. The V3 wing reaches L/D 8.6 in the wind tunnel (Poland et al. 2025) but 3.6-4.0 in traction with bridle and control unit (van der Vlugt et al. 2019), 3.9 fitted on a 160 m\textsuperscript{2} SkySails kite and 4-5 with tether (Erhard \& Strauch 2013, book chapter), 4-5 for SkySails towing kites (Paulig et al. 2013); 0.08 brings the aspect-ratio-3 kite with bridle to 4 at CL 0.8; 0 is a kite that keeps its design shape in flight. \\
rotor interference amplitude & \num{0.38} & \num{0}--\num{0.65} & -- & 3.9 & 0.0 & Loss amplitude of the rotor array factor 1 - (1 - 1/n) A exp(-s/s0), averaged over apparent wind angles of a reaching ship: a pair 3 diameters apart loses 4-8 \% in the wind tunnel at spin ratio 2 (Bordogna et al. 2020, Figs. 17-18), two to four rotors 9 diameters apart \textasciitilde{}1 \% (Han et al. 2025). The high bound is the fit with a free decay length. \\
kite landing max wave height & \num{4} & \num{3}--\num{5} & m & 3.7 & 0.0 & Estimate, no source gives a number: highest significant wave height for a landing that is not forced, the upper end of rough seas (sea state 5). The SkySails supplier reported limited launch opportunities in heavy sea (CE Delft 2016) without a number. \\
wing interference amplitude & \num{0.3} & \num{0.15}--\num{0.72} & -- & 3.6 & 0.0 & Loss amplitude of the wing array factor: 0.44 fits sails held at one angle of attack in 3D RANS (Malmek et al. 2024) and IDDES (Xu et al. 2025); independent sheeting recovers 0-80 \% of the loss depending on wind angle (van Reen 2026), a third taken here. \\
propeller diameter ratio & \num{0.5} & \num{0.45}--\num{0.65} & -- & 3.6 & 5.2 & Damen CF3850 2.6 m on 5.48 m draft (0.47), Ferus 4300 3.2 m on 5.95 m (0.54); electric drives allow larger, slower propellers, hence the high end. \\
wing pedestal width & \num{2.5} & \num{2}--\num{14} & m & 3.5 & 0.0 & Engineering estimate of the deck footprint of a wing raised on a pedestal above the container stacks; the high bound is a 14 m chord wing rooted at hatch level. \\
propeller viscous efficiency & \num{0.8} & \num{0.72}--\num{0.85} & -- & 3.4 & 4.6 & Engineering assumption: open-water efficiency is typically 0.55-0.70 (MAN ES 2018), which this factor reproduces on the actuator-disk efficiency; no direct source. \\
wing max lift coefficient & \num{2} & \num{1.6}--\num{2.5} & -- & 3.3 & 0.0 & 3D lift ceiling of a multi-element wing sail: AR 2 wings reach 1.70-1.73 with a 30 \% flap (Reche-Vilanova et al. 2021, Table 5), a single NACA0015 sail 1.08 (Xu et al. 2025, Table 4); BAR claims 2.5 times the lift of a single element for its three-element WindWing. Two-element sections reach 1.95-2.65 in 2D (Yao 2026). \\
canopy cloth area ratio & \num{3.5} & \num{3}--\num{4.3} & -- & 3.3 & 0.0 & Cloth sewn into a ram-air canopy per m\textsuperscript{2} of projected area: 3.0 m\textsuperscript{2} per flat m\textsuperscript{2} for two skins, ribs and diagonals (2.6-3.7), 1.17 flat m\textsuperscript{2} per projected m\textsuperscript{2}; paragliders weigh 0.19-0.21 kg per flat m\textsuperscript{2} with \textasciitilde{}35 g/m\textsuperscript{2} cloth (Advance Alpha 8, NOVA Prion and Mentor 2). \\
tether price & \num{93} & \num{40}--\num{300} & USD/kg & 3.1 & 0.0 & Purchase price of Dyneema rope at a series of 50 or more, EUR 80/kg (Joshi et al. 2025, p.703), paid to the rope maker without the market price factor; retail small diameters sell at 300-350 USD/kg (Marlow D12 via Mauripro/Binks); the low bound is an estimate for bulk rope of large diameter. \\
backhaul fraction & \num{0.3} & \num{0.15}--\num{0.5} & -- & 1.9 & 12.4 & Guadeloupe shipped 0.72 Mt out against 2.37 Mt in during 2024 (IEDOM annual report), including petroleum re-exports. \\
tether lifetime & \num{16000} & \num{4500}--\num{45000} & flight hours & 3.1 & 0.0 & Flight hours of a 12-strand SK78 tether flown at a fixed length; no ship or kite tether life has been measured, so this is our estimate from creep and its load peaks, the failure mode of a rope held under load without bending (Bosman et al. 2013; Humeau et al. 2018). Creep damage summed by Miner's rule over the hours flown at each tension on the routed trips of the 140 m, Kamsarmax and Neo-Panamax ships, at 25 °C: 10,000-12,100 h with the power law of Humeau et al. (exponent 4.43, 112 kJ/mol) anchored at 8 years at 20 \% of break load and 20 °C (DSM 2019), 22,700-26,100 h with the time-tension-temperature diagram of Bosman et al.; the base is their geometric mean, 15,100-17,800 h. The low bound is the power law at 30 °C with our harshest loads (4,700-5,500 h), the high bound the diagram at 20 °C (39,000-45,000 h). The swing of each figure eight, gusts and the ship's motion in a seaway shorten the creep life by 0-9 \% against a steady load at its mean, as the kite depowers at its force limit. Bending over the drum and sheaves at launch and landing is not in this life; it is charged per launch from the tether bend life. Bent over a drum under load in every pumping cycle, the rope would fail after 45 days, about 1,100 h (Bosman et al. 2013), and over a fixed fairlead in every figure eight within tens of hours (Davies et al. 2015), so the tow point turns with the tether; reeling in a few centimetres moves fresh rope against wear. Nearest service record: HMPE ship mooring lines, designed for 8,000 mooring hours to 75 \% strength at up to 22 \% of break load and retired by wear (Wilhelmsen 2022). The 7 years of the Ampyx Power fixed section (van Hagen et al. 2023) is an assumption with no stated basis. \\
rotor unit cost & \num{74000} & \num{36000}--\num{165000} & USD & 3.1 & 0.0 & Manufacturing cost of one rotor's fixed parts: least-squares line through bottom-up costs of Norsepower 2025 rotors at NREL composite and WISDEM steel costs; times the market price factor it reproduces Norsepower's 2019 prices to within -13 \% to +6 \%. \\
cargo value & \num{2500} & \num{1500}--\num{4000} & USD/t & 2.4 & 12.1 & Guadeloupe imported \textasciitilde{}EUR 3.9 bn on 2.37 Mt in 2024 (IEDOM annual report): \textasciitilde{}EUR 1,650/t over all imports and up to \textasciitilde{}EUR 3,000/t for non-bulk goods. \\
kite min elevation & \num{8} & \num{5}--\num{15} & deg & 2.8 & 0.0 & Engineering assumption: lowest tether elevation of a pattern, keeping the kite clear of the waves and the ship's superstructure. \\
service speed & \num{10.5} & \num{9.5}--\num{12} & kn & 2.8 & 2.2 & Conoship CIP3800 10 kn on 750 kW; the Damen CF3850's 11.8 kn is at 3 m ballast draft, so loaded speed is lower. \\
canopy cloth price per weight & \num{0.017} & \num{0.006}--\num{0.05} & USD/g & 2.7 & 0.0 & Rise of the goal price of a sewn m\textsuperscript{2} of cloth per g/m\textsuperscript{2} of cloth weight: the slope of the retail line, 0.031, scaled as canopy\_cloth\_base\_price; within Dacron sailcloth the price grows with weight to the power 0.6-1.0. \\
kite cycle loss probability & \num{5.0e-5} & \num{1.0e-5}--\num{0.005} & -- & 2.6 & 0.0 & Goal, half of about 1 in 10,000 per cycle for canopy or kite lost, shared with the tear. No source gives a rate; the high bound holds today's estimates. Chance per launch attempt of losing the kite in the water with its tether and control unit. \\
port time per call & \num{2} & \num{1}--\num{4} & days & 2.5 & 9.0 & Engineering assumption, independent of ship size: no current loading or discharge rate benchmark found (USDA tender terms of 2002-2011 give 5,000 t/day loading and 2,000 t/day discharging for bulkers). \\
rotor lift factor & \num{1} & \num{0.89}--\num{1} & -- & 2.5 & 0.0 & Factor on the lift of the ShipCLEAN rotor fit, CL 9.5 at spin ratio 3 corrected with full-scale measurements (Tillig \& Ringsberg 2020); CFD gives 8.5 (Li et al. 2012), the low bound. \\
kite launch min wind & \num{5} & \num{3}--\num{10} & m/s & 2.4 & 0.0 & Lowest apparent wind for a launch where it exceeds the wind derived from lift-off and flight. SkySails ships: about 5 m/s in perfect conditions, while rough sea and bad weather might require 10 m/s, the high bound (Paulig et al. 2013); not below Beaufort 3, 3.4-5.4 m/s (Konrad 2009). SkySails on land: 3 m/s (Fagiano et al. 2022), and 5 m/s measured against an assumed 6 (Junge et al. 2023). Kitepower Falcon: 5 m/s (Kitepower 2024). Rough sea comes with stronger wind, which clears the limit by itself, so the base is the calm-sea figure. \\
bridle cascade factor & \num{0.6} & \num{0.4}--\num{1} & -- & 2.4 & 0.0 & Engineering estimate: total bridle line length over the number of attachment lines times bridle height, from cascading lines into fewer main lines at 1:2 to 1:4 (Dunker 2013); 1 is no cascading. \\
rotor out of service & \num{0.03} & \num{0.01}--\num{0.15} & -- & 2.4 & 0.0 & Goal, the same for all rigs: share of round trips with a rotor out of service though the wind would allow it; units fail independently. Base: our reading of Estraden's operator, 97 \% usable over about 6 years, undefined (Spliethoff). Low: maker's above 99 \% on Estraden (Kuuskoski 2018). High: our estimate. \\
kite cycle damage probability & \num{5.0e-5} & \num{1.0e-5}--\num{0.01} & -- & 2.3 & 0.0 & Goal, half of about 1 in 10,000 per cycle for canopy or kite lost, shared with the loss at sea. No rate exists for a ship or large kite; the high bound holds today's estimates. Chance per launch attempt of a tear that ends the canopy. Gradual wear from flogging and folding is in the canopy life. Service limits of skydiving canopies, about 0.17 \% per jump, are retirement counts and only set the order of magnitude. \\
air density & \num{1.225} & \num{1.18}--\num{1.25} & kg/m\textsuperscript{3} & 2.2 & 0.2 & ISA sea level; humid tropical air is a few percent lighter. \\
block coefficient & \num{0.83} & \num{0.78}--\num{0.87} & -- & 2.1 & 0.9 & MAN ES (2018) lists coasters at C\_B 0.70-0.85; Watson-Gilfillan gives 0.79 at Fn 0.20; the dimensions and deadweight of the CF3850 and CIP3800 imply 0.85 or more. With 13.2 m beam it reproduces \textasciitilde{}3,800 dwt. \\
kite mass per area & \num{0.05} & \num{0.01}--\num{0.15} & t/m\textsuperscript{2} & 2.1 & 0.0 & Engineering assumption: a SkySails set-up of 160 m\textsuperscript{2} flat, 125 m\textsuperscript{2} projected, weighed \textasciitilde{}2 t (0.016 t/m\textsuperscript{2}, weak source), the kite alone under 1 t (EMSA 2023); 0.05 covers a heavier holding winch and mast. \\
holtrop max prismatic & \num{0.85} & \num{0.85}--\num{1} & -- & 2.0 & 5.2 & Largest prismatic coefficient at which the Holtrop-Mennen formulas are evaluated: the top of the tanker and bulk carrier range 0.73-0.85 and of the range of validity 0.55-0.85 tabulated by Ortigosa et al. (2009), second-hand, since Holtrop and Mennen (1982) and the pages of Holtrop (1984) we could read state no numbers; the 1982 paper limits the method to hulls resembling its average ship and calls its high-block adjustments tentative. 1 leaves the hull uncapped. \\
tether length ratio & \num{13} & \num{10}--\num{20} & -- & 2.0 & 0.0 & Tether length over projected span, held during flight and changed only at launch and landing (Erhard \& Strauch 2013). SkySails states a flat size for its 225 m\textsuperscript{2} kite (Paulig et al. 2013), and its 160 and 320 m\textsuperscript{2} kites are read the same way; at 160 m\textsuperscript{2} total to 125 m\textsuperscript{2} projected (Folkersma et al. 2020) they flew 300 m on 125 m\textsuperscript{2} (Erhard \& Strauch 2013, Fig. 18) and about 360 m on 250 m\textsuperscript{2} (Paulig et al. 2013). At the same projected area the aspect-ratio-3 span gives 15.5 and 13.1 spans; the base follows the larger kite. The high bound is the 19.7 spans of the 125 m\textsuperscript{2} kite's own 15.2 m span; the low bound is an estimate, a quarter shorter. \\
kite night launch & \num{1} & \num{0}--\num{1} & -- & 2.0 & 0.0 & Whether a kite is launched at night: 1 yes, 0 only between 06:00 and 18:00 local solar time. Operators said the SkySails kite could only be operated at daylight (CE Delft 2016); land kites have flown through the night (Fagiano et al. 2022, Junge et al. 2023). No source gives a night limit for a launch. \\
bridle height ratio & \num{0.7} & \num{0.5}--\num{0.9} & -- & 2.0 & 0.0 & Bridle length from canopy to control unit over projected span: parachutes rig A-lines of 0.5-0.6 constructed spans (Dunker 2013), 0.6-0.75 projected spans at a flat aspect ratio of 3.76 over 3; paraglider lines are \textasciitilde{}7-8 m on \textasciitilde{}11 m span. \\
kite max wind & \num{25} & \num{20}--\num{30} & m/s & 1.9 & 0.0 & SkySails ship kites were designed up to Beaufort 8; SkySails Power cuts out at 25 m/s at 200 m. \\
canopy mass overhead & \num{0.05} & \num{0.02}--\num{0.19} & kg/m\textsuperscript{2} & 1.9 & 0.0 & Canopy mass per m\textsuperscript{2} of projected area besides its cloth, at canopy\_cloth\_reference\_force and in proportion to the force limit, as they carry the load at the bridle attachments: seams, tapes, reinforcements and line attachments. Paragliders of 0.12-0.25 kg/m\textsuperscript{2} projected carry 0-20 g/m\textsuperscript{2} per flat m\textsuperscript{2} besides their cloth, lines and risers included (Ozone Ultralite 3 and 4, Nova Mentor 6, own fit); the high bound is what a 120 kg SkySails kite with lines, 225 m\textsuperscript{2} flat, 176 m\textsuperscript{2} projected, leaves beyond about 66 kg of cloth and 18 kg of lines, both own estimates (Paulig et al. 2013), which load patches for the \textasciitilde{}1 t point loads of a ship kite (Thedens 2015) could account for in part. \\
max froude number & \num{0.28} & \num{0.25}--\num{0.3} & -- & 1.9 & 0.1 & Holtrop-Mennen is validated up to about Fn 0.30 for general cargo ships; full coasters sail at Fn 0.23-0.29 (MAN ES 2018, coasters at 13-16 kn). \\
maintenance fraction & \num{0.03} & \num{0.02}--\num{0.05} & 1/year & 1.2 & 7.5 & Engineering assumption: maintenance of the hull and drive alone is commonly put at 2-5 \% of ship value per year; insurance comes on top, hence the high end. \\
wing unit cost & \num{66000} & \num{27000}--\num{153000} & USD & 1.9 & 0.0 & Engineering estimate of the manufacturing cost of a wing sail's controls, hydraulic power pack and 600 h of assembly; tilting or retraction is in the retractable rig premium. \\
rotor interference decay & \num{2.72} & \num{1.8}--\num{3.5} & diameters & 1.8 & 0.0 & Decay length of the rotor array loss, fitted to the wind-angle worst case of the Bordogna et al. (2020) pair at 3, 7.5 and 15 diameters; the low bound is the free fit with both parameters. \\
newbuild price exponent & \num{0.4} & \num{0.22}--\num{0.46} & -- & 1.4 & 7.3 & Exponent of price over deadweight from the same fit; 0.46 between 3,800 and 33,000 dwt, 0.22 between 33,000 and 82,000 dwt. Country of build differs between the small and large ships. \\
wing out of service & \num{0.03} & \num{0.004}--\num{0.15} & -- & 1.8 & 0.0 & Goal, the same for all rigs: share of round trips with a wing out of service though the wind would allow it; units fail independently. Base: our reading of the operator of the rotor ship Estraden, 97 \% usable, undefined (Spliethoff). Low: 99.6 \% operational availability on Canopee, maker with owner, undefined (OceanWings 2025). High: our estimate. \\
leeway limit & \num{10} & \num{6}--\num{15} & deg & 1.8 & 0.0 & Validity limit of the linear hull side-force model; captive tests on wind-assisted hulls cover 0-10 deg and ShipCLEAN reefs at 10 deg rudder angle (Tillig \& Ringsberg 2020). No operational limit found. \\
bridle rows & \num{4} & \num{3}--\num{8} & -- & 1.7 & 0.0 & Chordwise rows of line attachment points; ram-air canopies use A-D, and a large wing that keeps the same area per line needs A through H (Dunker 2013). \\
kite out of service & \num{0.03} & \num{0.01}--\num{0.15} & -- & 1.7 & 0.0 & Goal, the same for all rigs: share of round trips with the kite out of service though the wind and sea would allow it. No source gives a figure for any ship kite. Base and low: our assumption, the same as the rotors. High: our estimate. \\
vlsfo engine efficiency & \num{0.44} & \num{0.42}--\num{0.45} & -- & 1.7 & 2.1 & Wartsila 20: 190 g/kWh (0.444); MAN L21/31: 189-195 g/kWh (0.43-0.45) depending on attached pumps (project guide 2024). \\
wing endplate factor & \num{1.4} & \num{1.1}--\num{2} & -- & 1.6 & 0.0 & Effective over geometric aspect ratio from the deck mirror: 2.0 sealed, 1.36 at a root gap of 1 \% of span, 1.21 at 5 \%, 1.09 at 20 \% (minimum induced drag, Trefftz plane). WindWings stand about 7.5 m and the Oceanbird wing about 6 m above deck on pedestals (BAR; Oceanbird), the low bound. \\
tether drag coefficient & \num{1.1} & \num{1}--\num{1.2} & -- & 1.6 & 0.0 & Cylinder CD 1.1 at tether Reynolds numbers 2.7e3-8.2e3 (van der Vlugt et al. 2019, Table 1). \\
wing mass per area & \num{0.25} & \num{0.12}--\num{0.36} & t/m\textsuperscript{2} & 1.6 & 0.0 & Wing with pedestal and foundation: WindWing about 200 t for 750 m\textsuperscript{2} (DNV via Riviera 2024), Oceanbird Wing 560 about 200 t with foundation (Oceanbird); smaller hard sails 0.10-0.16 (EMSA 2023, Table 4; AYRO OW RE 45 t for 363 m\textsuperscript{2}). \\
deck length fraction & \num{0.75} & \num{0.65}--\num{0.8} & -- & 1.6 & 0.0 & Engineering assumption: share of the ship's length between forecastle and deckhouse that a row of rigs can use; an MR tanker has 140 m of upper deck on 182 m (Han et al. 2025), a CFD layout of five wing sails on a Kamsarmax spans 165 of 222 m (Mboumba 2022). \\
part load penalty & \num{0.2} & \num{0.1}--\num{0.3} & -- & 1.6 & 2.5 & Fitted to the MAN L21/31 SFOC table at 25-100 \% load: 0.14-0.27 depending on speed and attached pumps. \\
max air draft & \num{45} & \num{35}--\num{60} & m & 1.5 & 0.0 & Highest point of a deployed deck rig above the waterline; tilting or folding rigs pass under lower bridges. Panama new locks allow 57.9 m, the Kiel Canal \textasciitilde{}40 m. \\
rotor aspect ratio & \num{6} & \num{5}--\num{7} & -- & 1.5 & 0.0 & Height over diameter of the rotor the ShipCLEAN coefficients hold for (Tillig \& Ringsberg 2020); Norsepower builds 24x4, 30x5 and 35x5 m, Eco Flettner 18x3 m. \\
rotor installation fraction & \num{0.15} & \num{0.1}--\num{0.25} & -- & 1.4 & 0.0 & Installation over unit price, the foundation being in the unit price: 15 \% on a newbuild, 25 \% as a retrofit (EMSA 2023, Table 7). \\
reference draft & \num{5.5} & \num{5.25}--\num{5.95} & m & 1.4 & 4.8 & Conoship CIP3800 5.25 m, Damen CF3850 5.48 m, Ferus Smit 4300 5.95 m at summer draft (builder datasheets). \\
bridle spacing & \num{0.4} & \num{0.3}--\num{0.6} & -- & 1.3 & 0.0 & Engineering assumption: spanwise distance between attachment points over the mean chord, from paraglider line plans. \\
rotor mass per area & \num{0.63} & \num{0.3}--\num{0.7} & t/m\textsuperscript{2} & 1.3 & 0.0 & Eco Flettner 18x3 m weighs 34 t with a typical foundation (Interreg WASP, Annika Braren); a Norsepower 30x5 m is \textasciitilde{}42 t without foundation (EMSA 2023). \\
rotor lifetime & \num{25} & \num{20}--\num{30} & years & 1.3 & 0.0 & Anemoi states a 25-year design life; EMSA (2023) takes more than 20 years for wind-propulsion assets. \\
entrance length ratio & \num{0.22} & \num{0.15}--\num{0.3} & -- & 1.3 & 1.2 & Length of the bow entrance over waterline length; Holtrop's run length gives 0.26 L for the reference hull and the forward centre of buoyancy makes the entrance somewhat shorter. \\
availability & \num{0.92} & \num{0.88}--\num{0.95} & -- & 1.3 & 2.5 & Engineering assumption: dry-docking every 2.5 years plus repairs; no source found. \\
engine rating fraction & \num{0.9} & \num{0.75}--\num{0.9} & -- & 1.2 & 0.7 & Engine margin historically 10 \%, up to 25 \% (MAN ES 2018). \\
rotor max height & \num{35} & \num{30}--\num{40} & m & 1.2 & 0.0 & Height of the largest rotors built: 35x5 m tilting Norsepower rotors for the SC Connector (Marine Log 2020) and five 35x5 m Anemoi rotors on the Sohar Max (Maritime Executive 2024). \\
crew cost & \num{500000} & \num{350000}--\num{800000} & USD/year & 1.0 & 4.8 & Engineering assumption, no public source for coaster crew cost: \textasciitilde{}1,370 USD/day fits a mixed crew of seven, while 3,000 dwt coasters earned only \textasciitilde{}1,500-2,000 USD/day on time charter in 2024 (ISM), which implies low-cost crewing; an all-EU crew is at the high end. \\
launch mast max height & \num{35} & \num{29}--\num{60} & m & 1.2 & 0.0 & Tallest kite launch mast at sea, the 35 m Airseas mast; larger kites are flown from a mast of this height, beyond any built launch system. \\
kite stall margin & \num{2} & \num{1}--\num{4} & deg & 1.2 & 0.0 & Engineering assumption: margin below the angle of maximum lift that the bridle trim keeps in gusts. \\
pitch gyration ratio & \num{0.25} & \num{0.22}--\num{0.28} & -- & 1.2 & 1.8 & Pitch radius of gyration over length; 0.25 is the standard estimate for conventional cargo ships (ITTC). \\
crew cost exponent & \num{0.15} & \num{0}--\num{0.3} & -- & 1.0 & 4.6 & Exponent of crew cost over deadweight relative to the reference coaster: Handysize opex USD 4,500/day at 34.7k dwt and Supramax USD 4,910/day at 58.6k dwt (Pacific Basin 2024 interim report) minus maintenance\_fraction of their newbuild price leaves \textasciitilde{}USD 0.71-0.74 M/year of crew cost, against the 0.5 M coaster value. Drewry puts manning at over half of opex. \\
control unit mass per force & \num{0.22} & \num{0.15}--\num{0.45} & kg/kN & 1.1 & 0.0 & Mass of the airborne steering pod per kN of tether force limit: 70 kg for a 320 kN kite, steering forces up to 12 kN (Paulig et al. 2013), the one figure for a ship kite; its link and actuators carry loads in proportion to the force. Powered through the tether at 120 V, with a battery for emergencies (Maass \& Erhard 2013). The bounds are estimates. \\
launch mast steel cost & \num{2.51} & \num{1.3}--\num{5.2} & USD/kg & 1.1 & 0.0 & Manufacturing cost of fabricated tubular steel, 1.93 USD/kg (NREL WISDEM, Farkas \& Jarmai), with 30 \% added for the foundation. \\
midship coefficient & \num{0.99} & \num{0.98}--\num{0.995} & -- & 1.1 & 0.6 & Kerlen's regression gives 0.995 and Benford's 0.997 at C\_B 0.83. \\
ammonia tank gravimetric density & \num{0.5} & \num{0.33}--\num{0.7} & -- & 1.0 & 2.1 & Ammonia over the mass of the filled tank system. Our estimate: no source gives the mass of a ship's ammonia tank. \\
wing aspect ratio & \num{2} & \num{1.5}--\num{4} & -- & 1.0 & 0.0 & Installed wings span AR 1.9-3.8: WindWings 37.5 m x 20 m (BAR Technologies), Oceanbird Wing 560 40 m x 14 m, AYRO Oceanwings 363 m\textsuperscript{2} at 37 m. \\
lightship fraction & \num{0.28} & \num{0.25}--\num{0.32} & -- & 0.5 & 4.1 & Coasters carry 70-74 \% of their displacement as deadweight (Barlas, coaster main-dimension study); general cargo 65-80 \%. \\
launch mast mass coefficient & \num{0.00041} & \num{0.00027}--\num{0.00089} & t/m\^{}(7/3) & 1.0 & 0.0 & Engineering estimate of the steel mass of a telescopic launch mast, the base mass plus this coefficient times area\^{}(2/3) times height\^{}(5/3), for a tubular section that holds a launch pull of a tenth of the tether force; the tether runs to a deck tow point, not over the mast (Erhard \& Strauch 2013). \\
wing separation drag factor & \num{0.03} & \num{0.01}--\num{0.05} & -- & 1.0 & 0.0 & Fitted so that CD = 0.02 + 0.03 CL\textsuperscript{2} matches the 2D two-element section CD of 0.12 at CL 1.95 and 0.24 at CL 2.65 (Yao 2026). \\
ammonia pilot share & \num{0.05} & \num{0.05}--\num{0.15} & -- & 0.9 & 2.8 & Share of the fuel energy burned as VLSFO pilot in a dual-fuel ammonia engine: 5-15 \% expected, 5-10 \% on a diesel-cycle two-stroke (MMMCZCS 2021, Fuel Option Scenarios, p. 53). \\
discharge port dues & \num{0.1524} & \num{0.1524}--\num{0.6135} & EUR/m\textsuperscript{3} & 0.4 & 3.7 & Grand Port Maritime de la Guadeloupe, Droits de Port 2026 (Tarif 45), zone A, per m\textsuperscript{3} of length x beam x summer draught: general cargo ships of 65 m or more EUR 0.0964 inward plus 0.0560 outward; bulk carriers 0.3757 plus 0.2378. \\
rudder aspect ratio & \num{1.6} & \num{1.4}--\num{3} & -- & 0.9 & 0.1 & Effective aspect ratio of a rudder of geometric aspect ratio 1.5-2: CFD of a wind-powered cargo ship gives 0.93 times the geometric value, not the doubling of a mirror plane (Kramer \& Steen 2022); the high bound is that doubling. \\
launch mast telescope factor & \num{2.5} & \num{2.2}--\num{5} & -- & 0.9 & 0.0 & Price of the hydraulic telescopic launch mast over the steel cost of a static one: hydraulic marine cranes sell at 5.5-17 USD/kg against 1.0-1.6 USD/kg for steel monopoles, so the low bound is the cheapest crane per kg. The mast retracts because a kite aloft swings its tether round it; the SkySails ships used a telescopic mast, retracted in flight (Fritz 2013). \\
reference length & \num{85} & \num{84}--\num{88} & m & 0.4 & 3.5 & Real 3,800 dwt coasters: Damen Combi Freighter 3850 Lpp 84.97 m, Conoship CIP3800 Lpp 84.19 m; waterline length is \textasciitilde{}2-4 \% longer (MAN ES, Basic Principles of Ship Propulsion, 2018). \\
water viscosity & \num{1.19e-6} & \num{1.05e-6}--\num{1.35e-6} & m\textsuperscript{2}/s & 0.9 & 0.6 & Sea water at 15 °C (ITTC); tropical water is warmer and less viscous. \\
wing installation fraction & \num{0.15} & \num{0.1}--\num{0.25} & -- & 0.8 & 0.0 & Installation over unit price: 15 \% on a newbuild, 25 \% as a retrofit (EMSA 2023, Table 7). \\
kite forced landing loss probability & \num{0.05} & \num{0.01}--\num{0.2} & -- & 0.8 & 0.0 & Estimate, no source gives a rate: chance of losing the kite with its tether and control unit in a landing outside the landing limits, the only way down once wind or sea exceed the stay-up limits or at the end of a leg. \\
reference beam & \num{13.2} & \num{12.5}--\num{14} & m & 0.8 & 1.9 & Damen CF3850 12.5 m, Conoship CIP3800 13.2 m, Ferus Smit 4300 14.0 m (builder datasheets). \\
hull lifetime & \num{25} & \num{20}--\num{30} & years & 0.6 & 3.1 & Average ship demolition age has stayed above 27 years since 2019 and fell to 23 in the 2016 slump (Breakwave Advisors 2025). \\
source sea margin & \num{0.15} & \num{0.1}--\num{0.25} & -- & 0.8 & 2.2 & Sea margin inside the engine power of the class source designs, taken out to calibrate their calm-water resistance: 15 \% in every MAN ES table used (bulk carriers 2022 Tables 4-6, MR tankers 2019 Table 1, container vessels 2024 Table 19) and in the Knud E. Hansen PCTC consumptions; designers choose 10-30 \%, historically 15 \% (MAN ES 2023). \\
wing interference decay & \num{2.86} & \num{2}--\num{4} & chord lengths & 0.7 & 0.0 & Engineering estimate of the decay length of the wing array loss: no 3D data beyond 1.9 chords; set so that the loss of four wings is below 1 \% at 10 chords, the widest spacing in the 2D RANS of Wang et al. (2024). \\
wing lifetime & \num{25} & \num{20}--\num{30} & years & 0.7 & 0.0 & Oceanbird designs the Wing 560 for 25 years and BAR's builder claims 25 years; EMSA (2023) takes more than 20 years for wind-propulsion assets. \\
bulb area ratio & \num{0.06} & \num{0}--\num{0.1} & -- & 0.7 & 1.2 & Engineering assumption, no coaster data found; many coasters have plain bows, hence the low end of zero. \\
control unit cost per area & \num{15} & \num{8}--\num{30} & USD/m\textsuperscript{2} & 0.7 & 0.0 & Engineering estimate of the manufacturing cost of the three steering actuators, which grow with kite area: 3 x 5 kUSD at 1,000 m\textsuperscript{2}. \\
wave peakedness & \num{3.3} & \num{1}--\num{5} & -- & 0.7 & 1.2 & Peak enhancement factor of the JONSWAP spectrum: 3.3 is the mean of the JONSWAP measurements (Hasselmann et al. 1973); DNV-RP-C205 relates it to peak period and wave height, from 1 (Pierson-Moskowitz, fully developed sea) to 5 (steep wind sea). \\
deck mount height & \num{4} & \num{3}--\num{5} & m & 0.7 & 0.0 & Freeboard of real coasters is 1.5-2.0 m (CF3850, Ferus 4300, CIP3800 depth minus draft) plus a hatch coaming; no coaming height found. \\
ammonia engine premium & \num{190} & \num{90}--\num{300} & USD/kW & 0.7 & 0.9 & Dual-fuel ammonia two-stroke over a diesel: ME-LGI 590 against 400 USD/kW (Lloyd's Register \& UMAS 2020, Table B21, p. 60). Bounds by judgement. \\
retractable rig premium & \num{0.12} & \num{0.08}--\num{0.2} & -- & 0.6 & 0.0 & Extra price of tilting, folding or rails that lower wings and rotors for cargo handling and bridges, over the fixed rig price: bottom-up estimate of the hydraulics, hinge and foundation steel (10-17 \% for a tilting 35 x 5 m rotor of 56 t); such systems are in service (Norsepower tilting, Anemoi rails and folding, WindWings folding) but unpriced. \\
bulb height ratio & \num{0.4} & \num{0.3}--\num{0.6} & -- & 0.4 & 2.5 & Bulb area centre above keel over draft; Holtrop advises at most 0.6. \\
shaft efficiency & \num{0.96} & \num{0.95}--\num{0.97} & -- & 0.6 & 0.5 & Shaft with reduction gear 0.95-0.96 (MAN ES 2018). \\
kite launch max wind & \num{17} & \num{13}--\num{20} & m/s & 0.6 & 0.0 & Estimate, no source gives a launch limit: highest apparent wind for a launch, the top of Beaufort 7, one step below the Beaufort 8 flight limit of SkySails ship kites. SkySails on land reports landing as more time-critical than launch and reliable under nominal wind (Junge et al. 2023). \\
kite cycle course cost & \num{130} & \num{0}--\num{900} & USD & 0.5 & 0.0 & Estimate, no source gives a figure: mean cost per launch of turning for it. Lost sea time 0.1 h (0-0.7) times about 1,270 USD/h of hire and fuel. One hand launch needed a turn (WirtschaftsWoche 2008); no source says how many launches do. \\
empty container share & \num{0.188} & \num{0.16}--\num{0.188} & -- & 0.0 & 2.1 & Empty TEU over all TEU handled in EU27 main ports: 18.8 \% in 2024, 16.0-18.8 \% over 2010-2024 (Eurostat mar\_mg\_am\_cvh, loading status EMPTY and TOTAL). An empty box takes a slot and its tare but carries no cargo. \\
wing span efficiency & \num{0.9} & \num{0.8}--\num{0.95} & -- & 0.5 & 0.0 & Generic Oswald factor of a tapered wing; no wingsail-specific source found. \\
\end{longtable}
